\documentclass[11pt]{article}
\usepackage[margin=0.8in]{geometry}
\usepackage{amsmath,amssymb,amsthm}
\usepackage{enumitem}
\usepackage{longtable}
\usepackage{array}
\usepackage{xcolor}
\usepackage{hyperref}
\usepackage{fancyhdr}
\usepackage{titlesec}
\usepackage{tcolorbox}
\usepackage{multicol}
\usepackage{booktabs}
\hypersetup{colorlinks=true, linkcolor=blue!50!black, urlcolor=blue!50!black}
\pagestyle{fancy}
\fancyhf{}
\rhead{\thepage}
\lhead{Engineering Math to Pure Math}
\setlength{\parskip}{0.55em}
\setlength{\parindent}{0pt}
\setlist[itemize]{topsep=0.2em,itemsep=0.2em}
\setlist[enumerate]{topsep=0.2em,itemsep=0.25em}
\newtcolorbox{focusbox}{colback=blue!3!white,colframe=blue!50!black,title=Focus}
\newtcolorbox{intuitionbox}{colback=green!3!white,colframe=green!40!black,title=Intuition}
\newtcolorbox{formalbox}{colback=gray!4!white,colframe=gray!70!black,title=Formal}
\newtcolorbox{walkbox}{colback=orange!5!white,colframe=orange!60!black,title=Walking question}
\newcommand{\N}{\mathbb{N}}
\newcommand{\Z}{\mathbb{Z}}
\newcommand{\Q}{\mathbb{Q}}
\newcommand{\R}{\mathbb{R}}
\newcommand{\C}{\mathbb{C}}
\newcommand{\eps}{\varepsilon}
\newcommand{\pow}{\mathcal{P}}
\theoremstyle{definition}
\newtheorem{definition}{Definition}
\newtheorem{example}{Example}
\newtheorem{exercise}{Exercise}
\theoremstyle{plain}
\newtheorem{theorem}{Theorem}
\newtheorem{lemma}{Lemma}
\theoremstyle{remark}
\newtheorem{remark}{Remark}

\title{\textbf{From Engineering Mathematics to Rigorous Pure Mathematics}\\A Self-Study Curriculum for a Software Engineer}
\author{Prepared for Zim}
\date{May 2026}
\begin{document}
\maketitle
\tableofcontents
\newpage

\section{How to Use This Curriculum}

The first goal is not speed. The first goal is to become the kind of reader who can sit with a definition, test it against examples, see why each word is there, and then write a proof that another careful person can check.

The initial anchor text is Daniel Velleman's \emph{How to Prove It}. This curriculum prepares you to read it with less fog, then uses it as the first serious proof text. The material here is original and is meant to support the book, not replace it.

\begin{focusbox}
Each week has five parts: read, translate notation, prove small claims, debug one proof, and review. If a week feels too fast, split it into two. Mathematical maturity is built by returning to the same ideas with sharper eyes.
\end{focusbox}

\subsection{The Working Rhythm}

\begin{enumerate}
\item Read slowly with pencil in hand. Rewrite every definition in your own words.
\item Build a tiny example and a tiny non-example.
\item Try the proof before reading a polished solution.
\item Mark every proof gap with a concrete label: unknown definition, unjustified implication, hidden case, unclear quantifier, or algebra slip.
\item Review old prompts on days 1, 3, 7, 14, and 30.
\end{enumerate}

\section{The Map of Mathematics}

Engineering mathematics often asks, ``What method solves this problem?'' Pure mathematics asks, ``What structure is present, and what follows from the assumptions?''

\begin{longtable}{p{0.18\linewidth}p{0.37\linewidth}p{0.36\linewidth}}
\toprule
Area & Central objects & Why it matters later \\
\midrule
Logic & statements, quantifiers, implication & proof grammar for every later field \\
Sets and functions & membership, maps, images, inverse images & language of analysis, algebra, topology, geometry \\
Relations & equivalence, order, partitions & quotient spaces, congruences, bases, classification \\
Number systems & $\N,\Z,\Q,\R,\C$ & examples where proof methods become visible \\
Real analysis & limits, continuity, completeness & makes calculus honest and trains epsilon thinking \\
Abstract algebra & groups, rings, fields, modules & symmetry, equations, coding theory, cryptography, geometry \\
Topology & open sets, continuity, compactness & the portable skeleton of analysis and geometry \\
Differential geometry & manifolds, tangent spaces, curvature & calculus on curved spaces, physics, optimization \\
Algebraic geometry & varieties, ideals, schemes & shape described by polynomial equations \\
\bottomrule
\end{longtable}

\subsection{Prerequisite Dependency Graphs}

\begin{verbatim}
Logic -> Sets -> Functions -> Relations -> Proof fluency
Proof fluency -> Induction -> Countability -> Construction of number systems
Sets + Functions + Inequalities -> Sequences -> Limits -> Continuity
Limits + Completeness -> Real Analysis -> Metric Spaces -> Topology
Relations + Functions -> Groups -> Rings -> Fields -> Algebraic Geometry
Linear Algebra + Analysis -> Multivariable Analysis -> Differential Geometry
Topology + Commutative Algebra -> Algebraic Geometry
\end{verbatim}

\section{Notation and Logic Bootcamp}

\begin{definition}
A \emph{statement} is a sentence that is either true or false. A \emph{predicate} is a sentence with variables that becomes a statement after the variables are assigned values.
\end{definition}

\begin{intuitionbox}
Think of a predicate like a function returning a Boolean. In code, \texttt{isEven(n)} is not true or false until \texttt{n} is specified. In mathematics, $P(n)$ works the same way, except the domain and assumptions must be explicit.
\end{intuitionbox}

\begin{formalbox}
The statement $\forall n\in\Z,\ n^2\geq 0$ says every integer has a nonnegative square. The statement $\exists n\in\Z,\ n^2=2$ says at least one integer squares to $2$, which is false.
\end{formalbox}

\subsection{Core Symbols}

\begin{longtable}{p{0.16\linewidth}p{0.35\linewidth}p{0.38\linewidth}}
\toprule
Symbol & Meaning & Read it as \\
\midrule
$\forall$ & universal quantifier & for every \\
$\exists$ & existential quantifier & there exists \\
$\in$ & membership & is an element of \\
$\subseteq$ & subset & every element of the left set is in the right set \\
$\Rightarrow$ & implication & if left, then right \\
$\Leftrightarrow$ & biconditional & if and only if \\
$\neg P$ & negation & not $P$ \\
$P\land Q$ & conjunction & $P$ and $Q$ \\
$P\lor Q$ & disjunction & $P$ or $Q$, inclusive \\
$f:A\to B$ & function & $f$ maps each element of $A$ to an element of $B$ \\
$f^{-1}(S)$ & inverse image & all inputs whose outputs land in $S$ \\
\bottomrule
\end{longtable}

\subsection{Quantifier Negation}

\begin{theorem}
For any predicate $P(x)$ on a domain $D$,
\[
\neg(\forall x\in D,\ P(x)) \Longleftrightarrow \exists x\in D,\ \neg P(x).
\]
Also,
\[
\neg(\exists x\in D,\ P(x)) \Longleftrightarrow \forall x\in D,\ \neg P(x).
\]
\end{theorem}

\begin{intuitionbox}
To disprove ``every bridge in town is open,'' you need one closed bridge. To disprove ``some bridge is open,'' every bridge must be closed.
\end{intuitionbox}

\section{Weekly Study Plan: Bridge Into Velleman}

\begin{longtable}{p{0.05\linewidth}p{0.18\linewidth}p{0.25\linewidth}p{0.23\linewidth}p{0.19\linewidth}}
\toprule
Week & Theme & Main work & Exercises & Milestone \\
\midrule
1 & Statements and notation & Learn propositions, predicates, truth tables, domains & Translate 30 English statements into symbols and back & You can tell statement from predicate \\
2 & Quantifiers & Master $\forall,\exists$, nested quantifiers, negation & Write negations without the phrase ``it is not true that'' & Quantifier fog begins to lift \\
3 & Direct proof & Prove implications by assuming the hypothesis & Divisibility, parity, inequalities & First clean one-page proof \\
4 & Contrapositive and contradiction & Learn when each method reduces friction & Irrationality of $\sqrt2$, parity converses & You can choose a proof strategy \\
5 & Sets & Subsets, union, intersection, complement, power set & Prove set identities element by element & Set notation feels like grammar \\
6 & Functions & Domain, codomain, injective, surjective, bijective & Prove or refute claims about images and inverse images & You stop treating $f^{-1}$ as only an inverse function \\
7 & Relations & Equivalence relations, orders, partitions & Modular congruence, divisibility order & Quotients become less mysterious \\
8 & Induction & Weak, strong, structural induction & Sums, inequalities, recursive sequences & You can write induction without hand waving \\
9 & Velleman launch & Begin \emph{How to Prove It}, Chapters 1-2 & Redo early exercises in full sentences & Book reading becomes active \\
10 & Velleman logic & Conditional proof, biconditionals, cases & Build proof outlines before final proofs & You can see proof architecture \\
11 & Velleman sets & Set proofs, indexed families, power sets & Prove De Morgan laws for sets & Comfortable with element chasing \\
12 & Velleman functions & Composition, inverses, image/preimage & Prove composition preserves injectivity & Function proofs become routine \\
13 & Velleman relations & Equivalence classes and partial orders & Prove congruence mod $n$ is an equivalence relation & Quotient language opens \\
14 & Velleman induction & Strong induction and recursively defined objects & Fibonacci divisibility, binary strings & You can prove recursive facts \\
15 & Cardinality & Countable sets, diagonal arguments & Show $\Q$ is countable, $\R$ is uncountable & Infinity becomes structured \\
16 & Review project & Write a mini portfolio of 12 proofs & Revise old proofs with debugger rubric & Ready for analysis/algebra entry texts \\
\bottomrule
\end{longtable}

\section{Weekly Study Plan: After Velleman}

\begin{longtable}{p{0.05\linewidth}p{0.22\linewidth}p{0.32\linewidth}p{0.31\linewidth}}
\toprule
Week & Track & Study target & Milestone \\
\midrule
17 & Real analysis & Ordered fields, supremum property, absolute value & Prove basic inequalities from definitions \\
18 & Real analysis & Sequences and convergence & Write epsilon proofs for simple sequences \\
19 & Real analysis & Limits and continuity & Prove continuity using $\eps$-$\delta$ \\
20 & Real analysis & Compactness on $\R$ & Understand why closed and bounded matters \\
21 & Algebra & Groups, examples, subgroups & Prove a subset is or is not a subgroup \\
22 & Algebra & Homomorphisms and kernels & See kernels as algebraic fibers \\
23 & Algebra & Cosets and quotient groups & Prove Lagrange's theorem carefully \\
24 & Algebra & Rings, ideals, polynomial rings & Connect ideals to congruence \\
25 & Topology & Metric spaces and open balls & Rebuild continuity without formulas \\
26 & Topology & Open sets, closed sets, bases & Prove equivalent definitions of continuity \\
27 & Topology & Compactness and connectedness & Translate interval facts into topology \\
28 & Topology & Quotient spaces and product spaces & Recognize where algebra and topology rhyme \\
29 & Differential geometry & Curves, surfaces, charts & See a chart as a local coordinate API \\
30 & Differential geometry & Tangent vectors and differentials & Compute tangent maps and interpret them \\
31 & Differential geometry & Forms, integration, curvature preview & Connect local calculation to global shape \\
32 & Algebraic geometry & Affine varieties and polynomial ideals & Translate equations into geometric sets \\
33 & Algebraic geometry & Coordinate rings and maps & See functions on shapes as algebra \\
34 & Algebraic geometry & Hilbert basis theorem preview & Understand why finite generation is powerful \\
35 & Synthesis & Write a map of the fields & Explain how one theorem appears in three languages \\
36 & Portfolio & Ten polished proofs across fields & Ready for serious first courses \\
\bottomrule
\end{longtable}

\section{Theorem Summaries to Carry Forward}

\begin{longtable}{p{0.19\linewidth}p{0.34\linewidth}p{0.36\linewidth}}
\toprule
Theorem & Statement & What to notice \\
\midrule
Contrapositive & $P\Rightarrow Q$ is equivalent to $\neg Q\Rightarrow\neg P$ & Often turns a vague goal into concrete information \\
Euclid's Lemma & If $p$ is prime and $p\mid ab$, then $p\mid a$ or $p\mid b$ & The seed of unique factorization \\
Well-ordering of $\N$ & Every nonempty subset of $\N$ has a least element & The engine behind induction \\
Pigeonhole principle & More objects than boxes forces a collision & Counting can prove existence without construction \\
Cantor's theorem & No set has the same cardinality as its power set & Infinity has levels \\
Intermediate value theorem & Continuous functions on intervals hit intermediate values & Topology hiding inside calculus \\
Heine-Borel & In $\R^n$, compact means closed and bounded & A local-to-global machine for analysis \\
Lagrange's theorem & For finite group $G$, $|H|$ divides $|G|$ for every subgroup $H$ & Symmetry has arithmetic constraints \\
First isomorphism theorem & $G/\ker f\cong \operatorname{im} f$ & Quotient by what a map cannot see \\
Basis theorem preview & Polynomial ideals over a field are finitely generated & Algebraic geometry becomes finite enough to study \\
\bottomrule
\end{longtable}

\section{Formal Definitions With Friendly Pictures}

\subsection{Function}
\begin{formalbox}
A function $f:A\to B$ assigns to each $a\in A$ exactly one element $f(a)\in B$.
\end{formalbox}
\begin{intuitionbox}
The domain is not decoration. It is part of the machine. The formula $x^2$ behaves differently as a function $\R\to\R$, $\N\to\N$, or $\C\to\C$.
\end{intuitionbox}

\subsection{Injective and Surjective}
\begin{formalbox}
$f:A\to B$ is injective if $\forall a_1,a_2\in A,\ f(a_1)=f(a_2)\Rightarrow a_1=a_2$. It is surjective if $\forall b\in B,\exists a\in A$ such that $f(a)=b$.
\end{formalbox}
\begin{intuitionbox}
Injective means the function does not merge distinct inputs. Surjective means every target point is actually hit. In software terms, injectivity protects identity; surjectivity guarantees coverage of the declared output space.
\end{intuitionbox}

\subsection{Open Set}
\begin{formalbox}
In a metric space $(X,d)$, a set $U\subseteq X$ is open if for every $x\in U$, there exists $r>0$ such that $B_r(x)\subseteq U$.
\end{formalbox}
\begin{intuitionbox}
Every point of an open set has wiggle room. You can move a little without leaving the set. That small phrase becomes the foundation for continuity, manifolds, and topology.
\end{intuitionbox}

\section{Proof-Writing Drills}

\subsection{Daily Ten-Minute Drill}
\begin{enumerate}
\item Write the claim in symbols.
\item Circle the quantifiers.
\item Identify the hypothesis and conclusion.
\item List definitions that appear in the claim.
\item Decide the proof strategy: direct, contrapositive, contradiction, cases, induction, or construction.
\item Write a three-line proof skeleton.
\item Fill only one gap.
\item Read the proof aloud and remove any sentence that says ``clearly'' without evidence.
\end{enumerate}

\subsection{Walking Questions}

\begin{walkbox}
If a theorem says ``for every,'' what would a counterexample look like? Can you picture the single object that would break it?
\end{walkbox}

\begin{walkbox}
When a definition has three clauses, which clause is easiest to forget? What example satisfies two clauses but fails the third?
\end{walkbox}

\begin{walkbox}
What is the difference between proving $A\subseteq B$ and proving $A=B$? Where do the two directions live?
\end{walkbox}

\begin{walkbox}
What does a function remember, and what does it forget? In algebra, that question becomes the idea of a kernel.
\end{walkbox}

\section{Proof Debugger Workflow}

When you submit or review a proof, use this structure.

\begin{enumerate}
\item \textbf{Claim parsed:} Rewrite the theorem with all quantifiers explicit.
\item \textbf{Definitions needed:} List every definition used.
\item \textbf{Strategy:} Name the method and why it fits.
\item \textbf{Line check:} For each line, ask what justifies it.
\item \textbf{Gap labels:} Mark any issue as quantifier gap, definition gap, implication gap, case gap, example gap, or notation gap.
\item \textbf{Minimal repair:} Fix the smallest broken step before rewriting the whole proof.
\item \textbf{Final polish:} Remove computational clutter and keep the logical spine visible.
\end{enumerate}

\begin{focusbox}
The best feedback is specific. ``This proof is wrong'' is smoke. ``Line 4 assumes the conclusion for an arbitrary $x$, but $x$ was chosen from a smaller set'' is fire you can put out.
\end{focusbox}

\section{Milestones Toward Higher Mathematics}

\subsection{Real Analysis}
You are ready when you can negate nested quantifiers, prove inequalities from definitions, and write epsilon proofs without guessing the epsilon at the end. Analysis is where the infinite becomes accountable.

\subsection{Abstract Algebra}
You are ready when functions, equivalence relations, and modular arithmetic feel natural. Algebra studies operations and symmetry. A group is not a bag of elements; it is a world where composition behaves predictably.

\subsection{Topology}
You are ready when open sets, inverse images, and proof by definition feel comfortable. Topology studies what survives stretching, bending, and continuous deformation. It is the part of analysis that stopped depending on distance.

\subsection{Differential Geometry}
You are ready when linear algebra, multivariable calculus, and topology have become friends. A manifold is a space that may be globally curved or complicated but locally looks like $\R^n$.

\subsection{Algebraic Geometry}
You are ready when rings, ideals, quotient structures, and topology no longer sound like separate countries. Algebraic geometry studies geometric shapes by the polynomial equations that define them.

\section{Historical Thread}

Proof culture did not appear all at once. Euclid organized geometry into axioms and deductions. Fermat, Euler, Gauss, Cauchy, Weierstrass, Noether, Hilbert, Grothendieck, and many others changed what mathematicians considered a satisfying explanation. Calculus began as a powerful engine for motion and change; analysis later rebuilt its foundations. Algebra began with solving equations; it became the study of structure. Geometry began with measurement; it became the study of spaces, transformations, and local coordinates.

This matters because the fields are not arbitrary school subjects. They are repairs and expansions. When a method works too well to ignore but too vaguely to trust, mathematics often responds by inventing a sharper language.

\section{What to Focus On}

\begin{itemize}
\item Quantifiers are the steering wheel.
\item Definitions are not vocabulary; they are tools with exact operating instructions.
\item Examples prevent you from floating away.
\item Counterexamples teach the boundary of a theorem.
\item A proof is a chain of justified transitions, not a performance of confidence.
\item The same ideas recur: structure, maps, invariants, quotients, limits, compactness, symmetry.
\end{itemize}

\end{document}
